%% ============================================================
%%  4.  METHOD: ADAPTIVE INFLUENCE BALANCING
%% ============================================================
\section{Adaptive Influence Balancing}
\label{sec:method}

We introduce \textbf{Adaptive Influence Balancing} (AIB), a framework that
monitors live process diagnostics during multi-agent debates and selectively
intervenes to protect minority agents whose reasoning shows signs of being
prematurely eliminated.

%% ----------------------------------------------------------
\subsection{Problem Formulation}
\label{sec:formulation}
%% ----------------------------------------------------------

Let $N$ agents debate a multiple-choice question $x$ with label set
$\mathcal{H}$ over $T$ rounds.
After round $t$, agent $i$ holds answer $h_i^{(t)}\!\in\!\mathcal{H}$,
confidence $c_i^{(t)}\!\in\![0,1]$, and a free-text explanation
$e_i^{(t)}$.
The \emph{support share} of hypothesis $h$ at round $t$ is
\begin{equation}
  s_t(h)=\frac{1}{N}\sum_{i=1}^{N}\mathbf{1}[h_i^{(t)}=h].
\end{equation}
Agent $i$ is a \emph{minority agent} when its current answer has strictly
less support than the most popular answer:
$s_t(h_i^{(t)}) < \max_{h\in\mathcal{H}} s_t(h)$.

We introduce two evaluation metrics beyond accuracy.
\textbf{Consensus Reliability} (CR) measures whether reached consensus is
correct:
$\text{CR} = P(\text{correct}\mid\text{consensus reached})$.
\textbf{Premature Elimination Rate} (PER) measures the fraction of correct
hypotheses discarded before the final round:
\begin{equation}
  \text{PER} = \frac{\bigl|\{h : h=y,\;h\!\in\!\mathcal{H}_{\mathrm{proposed}},\;h\!\notin\!\mathcal{H}_{\mathrm{final}}\}\bigr|}
                    {\bigl|\{h : h=y,\;h\!\in\!\mathcal{H}_{\mathrm{proposed}}\}\bigr|},
\end{equation}
where $y$ is the ground-truth answer, $\mathcal{H}_{\mathrm{proposed}}$ is
the set of correct answers ever proposed, and $\mathcal{H}_{\mathrm{final}}$
is the set of correct answers still held at the last round.
A high PER indicates that the debate discards valuable reasoning before it
can influence the final answer.

%% ----------------------------------------------------------
\subsection{Intervention Mechanisms}
\label{sec:mechanisms}
%% ----------------------------------------------------------

We design three complementary mechanisms.
Each is evaluated as a standalone condition; they share the same base
debate protocol and differ only in the trigger logic and the additional
prompt context given to flagged agents.

\paragraph{Loser Resistance (Exp2a/2b).}
Agents that lose peer support are prompted to re-examine their position
from first principles rather than following the majority.
A \emph{resistance scalar} $\rho_i^{(t)}\!\in\![0,\rho_{\max}]$ is
maintained per agent; it increases by $\alpha\!=\!0.08$ per round that
agent $i$ is in the minority, capped at $\rho_{\max}\!=\!0.65$.

\textbf{Exp2a} (support-triggered): agent $i$ is flagged when its support
share falls below 1.5 agents and the trend over the last two rounds is
negative.
\textbf{Exp2b} (metric-triggered): replaces the raw-support heuristic with
live process diagnostics—influence asymmetry $\hat{I}>0.65$ or balance
$B<0.30$—providing a principled, interpretable trigger.

A flagged agent receives a prompt suffix instructing it to switch
\emph{only if} a peer argument concretely defeats its reasoning; otherwise
it should preserve its answer and explain why.

\paragraph{Static Minority Protection (Exp3a).}
After Round~1, any agent whose answer differs from the two-agent majority is
flagged as a minority and given two protections: (i) an extra $+110$ token
budget for its next response, and (ii) an instruction to first generate the
strongest evidence for its hypothesis and then the strongest evidence against
it before deciding whether to update.
Exp3a uses a purely categorical trigger: one agent differs from the other two.

\paragraph{Diagnostic-Triggered Minority Protection (Exp3b).}
Exp3b replaces the categorical trigger with live process diagnostics.
An agent is flagged when \emph{either}:
\begin{itemize}[noitemsep,topsep=2pt]
  \item influence asymmetry $\hat{I}\!\geq\!0.65$
        \textbf{and} the agent's influence share $<\!0.34$, or
  \item balance $B\!\leq\!0.30$ (rapid convergence detected),
\end{itemize}
where $\hat{I}$ and $B$ are computed from the trajectories accumulated so
far following the formulas of \citeauthor{diagnostic2026}~(\citeyear{diagnostic2026}).
Flagged agents receive the same $+110$ token budget and protect-before-change
instruction as Exp3a.
The diagnostic trigger is more selective than the categorical one: it
activates only when the debate dynamics are actually imbalanced, not
whenever any disagreement exists.

%% ----------------------------------------------------------
\subsection{Learned Predictor (Exp4)}
\label{sec:predictor}
%% ----------------------------------------------------------

The hard thresholds in Exp2b and Exp3b were set by inspection.
In Exp4 we replace them with a \emph{learned predictor} that estimates the
probability that a minority hypothesis is correct given the current debate
state.

\paragraph{Training data.}
We extract every minority situation from the 200 baseline debates (Exp1):
for each question at each of rounds 2--4, and for each answer with strictly
less than maximum support, we record the current feature vector and label
$z=\mathbf{1}[h_{\mathrm{minority}}=y]$.
This yields approximately 840 labeled minority situations with a positive
rate of $\approx\!18\%$ (correct minorities are rare but important).

\paragraph{Feature engineering.}
We combine three signal groups:
\begin{enumerate}[noitemsep,topsep=2pt]
  \item \emph{Process diagnostics}: influence asymmetry $\hat{I}$, balance
        $B$, engagement, responsiveness (computed live up to the current round).
  \item \emph{Minority-specific}: support share $s_t(h)$, support delta
        $\Delta s_t(h)$, mean LLM-judged explanation quality
        $\bar{q}_{\mathrm{min}}$, mean confidence $\bar{c}_{\mathrm{min}}$,
        number of defections from $h$, and whether $h$ was the majority in
        the previous round.
  \item \emph{Context}: rounds remaining ($T-t$), number of current supporters.
\end{enumerate}

\paragraph{Model.}
We train a logistic regression with $\ell_2$ regularization selected via
5-fold cross-validation, followed by isotonic-regression calibration on a
held-out validation split~\cite{pedregosa2011scikit}.
The top predictors are minority reasoning quality ($\beta=0.245$), support
trajectory ($\beta=0.179$), and defection count ($\beta=-0.121$), validating
that the LLM judge score is causally predictive of correctness.

\paragraph{Adaptive token budget.}
Instead of a fixed $+110$ tokens, Exp4 scales the extra budget by the
predicted probability:
\begin{equation}
  \Delta_{\mathrm{tokens}}(i,t)
    = \bigl\lfloor 150 \cdot \min\!\bigl(p_i^{(t)},\;0.9\bigr) \bigr\rfloor,
\end{equation}
where $p_i^{(t)} = P(h_i^{(t)}=y \mid \phi_i^{(t)})$ is the calibrated
predictor output.
An agent is protected only when $p_i^{(t)}\!\geq\!0.35$ and
$\bar{q}_{\mathrm{min}}\!\geq\!0.30$, ensuring that quality-less dissent
receives no extra resources.
